\section{Diálogo entre \enquote{Ensaio sobre a Dádiva} (Mauss) e \enquote{Reciprocidade e circulação de crianças entre camponeses no sertão} (Petrafesa de Godoi)}

A prática de se pegar crianças de pais \enquote{mais fracos} me chamou atenção para algo que se assemelha à compreensão de Mauss sobre a esmola;
    primeiramente, me pareceu haver ali este aspecto de sacrifício feito ao sagrado -- de um \textit{potlatch} oferecido aos deuses e espíritos de antepassados -- dado que:
        segundo Mauss, estes \enquote{consentem que as porções que lhes dão e que são destruídas em sacrifícios inúteis sirvam aos pobres e às crianças} (Mauss, p. 207), e o aceite de uma criança ofertada consiste num compromisso de queima de riqueza, uma promessa de \textit{potlatch} contínuo a partir do momento do aceite.

Receber o filho de um casal \enquote{mais fraco} equivaleria a acender uma fogueira de \textit{potlatch} e mantê-la acesa, alimentando seu fogo, por quanto tempo se prolongar a criação da criança recebida;
    o caráter agonístico aparece aí não em relação aos doadores, mas aos demais membros da comunidade:
        os casos de três dos donatários citados -- seu Manuel Lourenço, dona Maria Coelho, e dona Ana -- ajudam a ilustrar este movimento.

Quando sua filha morre, seu Manuel assume a criação de um dos sete filhos deixados por ela e repartidos entre parentes;
    e o artigo destaca uma \enquote{honra do pai de família} como motor para o compromisso de seu Manuel:
        \enquote{Ajuda} é o termo utilizado pelos camponeses entrevistados para se referir ao recebimento de crianças de outras famílias.
        
Esta nomeação parece indicar que receber uma criança funciona, de um lado, como o recebimento de uma dádiva dada pela família de origem, mas, de outro, ela é a própria dádiva, que parte da família que a recebe;
mais precisamente, ela é um sacrifício, feito não para a família de origem, mas para uma comunidade que assume a posição de um sagrado.
O caso de seu Manuel é paradigmático porque a família de origem da criança não existe, ela está órfã, então seu neto não assume o papel de dádiva ofertada, a única dádiva que está posta é a da \enquote{ajuda}, um sacrifício necessário para que seu Manuel mantenha sua \enquote{honra de pai de família} perante sua comunidade;
a \enquote{ajuda} é um \textit{potlatch} que ele oferta à sua comunidade.

O caso de dona Maria é interessante pois a relação aassimétrica entre ela e a damília de origem de Leciana, a criança assumida por ela, lhe permite devolver a dádiva recebida sem constrangimento social.
Aqui não aparece um caráter sacrificial, a \enquote{ajuda} está condicionada ao cumprimento de certas funções pela criança recebida.
A assimetria me faz pensar numa relação entre uma chefia -- dona Maria -- e um agregado chefiado -- a família de origema de Leciana --, que oferece a filha como uma dádiva a fim de seduzir a chefia a um laço mais íntimo e duradouro.
Mas dada a posição relativa superior de dona Maria, Leciana não tem um papel de tonga -- ao contrário do caso do neto de seu Manuel, uma tonga de sua própria família que seu Manuel vai resgatar para manter sua honra --, mas está muito mais para a de oloa, pois funciona como um instrumento de trabalho para dona Maria.
O laço íntimo estabelecido com a família de Leciana passa a lhe desfavorecer quando Leciana engravida solteira na adolescência e, não havendo uma questão de honra a se manter perante a comunidade, Leciana é devolvida sem prejuízos simbólicos para dona Maria, o que não aconteceria com seu Manuel caso negasse assumir a criação de seu neto.

Mas o mais paradigmático de todos os casos é o de dona Ana, que assume nove crianças além de seus próprios onze filhos.
As crianças vinham tanto de parentes quanto de desconhecidos.
Chama atenção que ela mesma chegou a pedir pela criança de outras pessoas, o que fortalece um papel sacrificial da \enquote{ajuda} aqui.
Trata-se de uma situação complementar à de seu Manuel, que se viu forçado a realizar a \enquote{ajuda} para manter sua honra;
ao pedir pela criança, dona Maria oferece um sacrifício que não seria, ao menos a princípio, obrigada a fazer em nome de sua honra.
No entanto, sua \enquote{ajuda} tem função retributiva e sacrificial.
É onde aparece mais evidente uma forma de \textit{potlatch}, pois a quantidade de crianças, 20 no total, criadas por dona Ana, revela uma vida de riquezas adquiridas por ela numa relação contínua de trocas com sua comunidade, riquezas que se demandariam de um chefe que sejam distribuídas ou queimadas.
Cada \enquote{ajuda} de dona Ana funciona como uma fogueira \textit{potlatch}, e o acúmulo dessas fogueiras faz com que sua fumaça comece a alcançar territórios cada vez mais distantes, elevando sua honra na medida que seu nome vai sendo conhecido comoo de uma pessoa generosa e \enquote{forte}, conforme indica o artigo.
Esta honra proporciona a dona Ana mais oportunidades de ofertar \enquote{ajuda}, pois pessoas de povoados vizinhos a procuram para lhe entregar seus filhos, e quanto mais fogueiras são acumuladas, mais \enquote{forte} fica seu nome, o que atrai propostas de alianças diversas por pessoas hierarquivamente innferiores a ela em sua comunidade, melhorando suas condições de prosperidade material, de aquisição de oloas.

Mas há peculiaridades muito interessantes no caso da \enquote{ajuda}, pois o sistema de prestações totais parece estar em curto-circuito.
A \enquote{ajuda} é um ato contínuo, não um evento como o \textit{potlatch} tradicional.
A queima de riqueza se dá continuamente, assim como essa fogueira sacrificial funciona como tonga num laço que transcende a família de origem e compreende a comunidade como um todo;
é um sacrifício em ato contínuo à comunidade enquanto sagrado, que funciona como canal para o estabelecimento de novos laços, ou seja, como um canal para novos canais, para uma crescente aquisição tanto de oloas quanto de tongas.
Este funcionamento da \enquote{ajuda}manifesto em maior evidência pelo caso de dona Ana apresenta grande semelhança de forma com o do capital nas sociedades capitalistas, dado que este se estabelece como canal para novos canais apenas enquanto seu detentor lhe sustenta queimando suas riquezas, isto é, investindo em ato contínuo.

A \enquote{ajuda} parece ainda um passo no sentido de infividualizar a pessoa moral, completado pelo capital.
O sacrifício ao sagrado é uma relação distinta daquela do \textit{potlatch} tradicional, pois, de um lado, pode oferecer o sacrifício tanto a pessoa individual quanto um grupo familiar, um clã, etc, estando o sagrado numa posição hierárquica absolutamente superior, não se tratando, portanto, de uma disputa de honra entre pessoas morais.
Trata-se, na veradde, de uma tentativa de seduzir o sagrado, de adquirir hona aos seus olhos, prestando-se a retribuir pelas riquezas recebidas em troca da oportunidade de realizar novas trocas e receber mais riquezas.
pE o que faz a família de Leciana ao oferecê-la para dona Maria:
oferece a filha em sacrifício em prol de um laço mais forte com uma pessoa hierarquicamente superior.
Acontece que dona Maria, no papel de sagrado, aceita apenas temporariamente o sacrifício, rejeitando-o quando este não mais lhe agrada.
As crianças aqui podem funcionar tanto como oloas quanto como tongas, ao mesmo tempo, para pessoas distintas ou para a mesma pessoa.
Há uma sobreposição de papéis assim como há uma sobreposição de \textit{potlatchs} acontecendo em níveis distintos.
Seu Manuel tem o neto como tonga, o que parece até redundante se dizer.
Leciana é tonga para sua família de origem, mas apenas oloa para dona Maria, pois sua \textit{ajuda} não tem caráter sacrificial, e a posse de Leciana não estabelece um canal substantivo para oloas além do trabalho doméstico pela própria Leciana.
Dona Ana, por fim, ainda que possa vir a usufruir do trabalho das crianças que assume, o que as colocaria no papel de oloa, claramente as possui como tonga diante do sagrado a que oferece este, fortalecendo um canal pelo qual recebe tanto oloas quanto tongas.
A \enquote{ajuda} não é um sacrifício de toda a comunidade, mas de seu Manuel, dona Maria (em menor grau) e dona Ana (em maior grau), ou, no máximo, dos núcleos familiares de cada um deles.
A comunidade, que nomeia o recebimento de crianas como \enquote{ajuda}, é a pessoa moral no papel de sagrado à qual cada um oferece a \enquote{ajuda} em sacrifício. O capital individualiza ainda mais esta relação, pois o indivíduo é que possui o capital (tonga), mas apenas na medida em que oferece dinheiro (oloa) em sacrifício a um sagrado -- a sociedade capitalista -- na forma de investimentos, buscando igualmente fortalecer um canal para o recebimento de mais tongas e oloas.